% Zone „Das kennst du schon“ – aus bank/schnittmengen/zone.jsonl (zusammenbau v0.1)
\einheitenkopf[zone][Kennst du schon]{Das kennst du schon}
\setcounter{aufgabe}{0}

%% TODO Titel: Ich-kann-Satz fehlt in der Bank (Platzhalter: Kettenname) – Nr. 1
%% TODO Anweisung: ein Satz über der ersten Teilaufgabe fehlt in der Bank – Nr. 1
\begin{aufgabe}{Geradengleichungen aufstellen und den allgemeinen Geradenpunkt bilden}
\begin{teile}
\teil Gerade durch $A(1|2|0)$ und $B(3|5|1)$ – Gleichung in Parameterform? x = \leerfeld + t · \leerfeld
\teil $\vec{x} = (3|-1|0{,}5) + t \cdot (-2|4|1{,}5)$ – Punkt für $t = -0{,}5$? ( \leerfeld | \leerfeld | \leerfeld )
\end{teile}
\end{aufgabe}

%% TODO Titel: Ich-kann-Satz fehlt in der Bank (Platzhalter: Kettenname) – Nr. 2
%% TODO Anweisung: ein Satz über der ersten Teilaufgabe fehlt in der Bank – Nr. 2
\begin{aufgabe}{Ebenengleichungen lesen (Koordinatenform, Parameterform, Normalenvektor)}
\begin{teile}
\teil $2x - y + 3z = 6$ – Normalenvektor? ( \leerfeld | \leerfeld | \leerfeld )
\teil Liegt $Q(-4|7|1)$ in der Ebene $0{,}5x + 4z = 2$? ☐ ja ☐ nein
\end{teile}
\end{aufgabe}

%% TODO Titel: Ich-kann-Satz fehlt in der Bank (Platzhalter: Kettenname) – Nr. 3
%% TODO Anweisung: ein Satz über der ersten Teilaufgabe fehlt in der Bank – Nr. 3
\begin{aufgabe}{Den Lagebefund vorweg führen (liegt der Punkt darin, schneidet die Gerade überhaupt – Merkregel über das Skalarprodukt)}
\begin{teile}
\teil $(1|2|1) \cdot (2|-1|0)$ – Skalarprodukt? \leerfeld
\teil Ist die Gerade $\vec{x} = (1|0|2) + t \cdot (2|1|-1)$ parallel zur Ebene $x - 2y + 3z = 5$? ☐ ja ☐ nein
\end{teile}
\end{aufgabe}

%% TODO Titel: Ich-kann-Satz fehlt in der Bank (Platzhalter: Kettenname) – Nr. 4
%% TODO Anweisung: ein Satz über der ersten Teilaufgabe fehlt in der Bank – Nr. 4
\begin{aufgabe}{Lineare Gleichungssysteme mit zwei und drei Unbekannten lösen, auch überbestimmte mit Probegleichung}
\begin{gleichungsraster}[2]
\gl{\text{$r + s = 5$ und $r - s = 1$ – $r$ und $s$?}} &
\gl{\text{$2r + s = 1$ und $r - s = 3{,}5$ – $r$ und $s$?}} \\
\end{gleichungsraster}
\end{aufgabe}

%% TODO Titel: Ich-kann-Satz fehlt in der Bank (Platzhalter: Kettenname) – Nr. 5
%% TODO Anweisung: ein Satz über der ersten Teilaufgabe fehlt in der Bank – Nr. 5
\begin{aufgabe}{Teilverhältnisse und Parameterbereiche von Strecken lesen}
\begin{teile}
\teil Die Strecke $AB$ ist $\vec{x} = (0|0|0) + t \cdot (4|2|6)$ mit $0 \le t \le 1$. Punkt für $t = 0{,}5$? ( \leerfeld | \leerfeld | \leerfeld )
\teil Die Strecke $AB$ ist $\vec{x} = (2|-1|0) + t \cdot (6|3|3)$ mit $0 \le t \le 1$. Punkt $T$ für $t = \frac{1}{3}$ und Teilverhältnis $AT : TB$? T( \leerfeld | \leerfeld | \leerfeld ), AT : TB = \leerfeld : \leerfeld
\end{teile}
\end{aufgabe}

%% TODO Titel: Ich-kann-Satz fehlt in der Bank (Platzhalter: Kettenname) – Nr. 6
%% TODO Anweisung: ein Satz über der ersten Teilaufgabe fehlt in der Bank – Nr. 6
\begin{aufgabe}{Schrägbilder lesen und zeichnen}
\begin{teile}
\teil Zeichne den Punkt $P(2|3|1)$ in das Schrägbild ein.

\begin{ksys3}[xyz,x1min=0,x1max=4,x2min=0,x2max=4,x3min=0,x3max=3] \end{ksys3}
\teil Zeichne den Punkt $R(-1|2|1{,}5)$ in das Schrägbild ein.

\begin{ksys3}[xyz,x1min=-2,x1max=3,x2min=0,x2max=4,x3min=0,x3max=3] \end{ksys3}
\end{teile}
\end{aufgabe}

%% TODO Titel: Ich-kann-Satz fehlt in der Bank (Platzhalter: Kettenname) – Nr. 7
%% TODO Anweisung: ein Satz über der ersten Teilaufgabe fehlt in der Bank – Nr. 7
\begin{aufgabe}{Geradengleichungen aufstellen und den allgemeinen Geradenpunkt bilden – Fehler finden}
\begin{teile}
\teil Ich finde den Fehler: Zu $\vec{x} = (2|1|0) + t \cdot (1|3|-1)$ schreibt Tim den allgemeinen Geradenpunkt als $(t|3t|-t)$.
\end{teile}
\end{aufgabe}

%% TODO Titel: Ich-kann-Satz fehlt in der Bank (Platzhalter: Kettenname) – Nr. 8
%% TODO Anweisung: ein Satz über der ersten Teilaufgabe fehlt in der Bank – Nr. 8
\begin{aufgabe}{Geradengleichungen aufstellen und den allgemeinen Geradenpunkt bilden – selbst rechnen}
\begin{teile}
\teil $\vec{x} = (-1|3|2) + t \cdot (2|0|1)$ – allgemeiner Geradenpunkt? ( \leerfeld | \leerfeld | \leerfeld )
\end{teile}
\end{aufgabe}
